% Faithful conversion of chapter-05-text-h-0097; source ODF IDs and text order retained.
\hypertarget{__RefHeading___Toc4825_3118176633}{}
\refstepcounter{section}\section*{4.4 Conclusão}
\label{ages-iii:text-h-0097}
\addcontentsline{toc}{section}{4.4 Conclusão}
% text-p-1310
\indent \originaltab{}Ao longo desta edição da AGES, tive a oportunidade de participar de um projeto bastante diferente dos que costumam ser desenvolvidos na disciplina. A integração de uma aplicação de software com um hardware proprietário, trouxe desafios técnicos que vão além do desenvolvimento tradicional de aplicações web. Diferentemente de projetos focados apenas em um frontend, backend e banco de dados, este exigiu que toda a equipe compreendesse o funcionamento do leitor RFID, as limitações do equipamento e a forma como essas informações seriam transportadas até a aplicação. Esse cenário tornou o desenvolvimento mais complexo, mas também muito mais enriquecedor.\par

% text-p-1311
\indent A complexidade do projeto impactou todos os níveis da equipe. Os AGES I tiveram o desafio de desenvolver funcionalidades em um contexto técnico bastante exigente; os AGES II precisaram modelar estruturas de dados enquanto muitos requisitos ainda estavam sendo descobertos; os AGES III foram responsáveis por transformar essas necessidades em uma arquitetura consistente, integrando infraestrutura em nuvem, backend, frontend e hardware; e os AGES IV tiveram o desafio de conduzir um projeto cujo domínio de negócio também era novo para eles. Apesar das dificuldades, o comprometimento da equipe permitiu que o projeto evoluísse de forma consistente até sua conclusão.\par

% text-p-1312
\indent Para mim, esta experiência também representou um importante marco de evolução profissional. Minha última participação na AGES, como AGES II, havia ocorrido há cerca de dois anos. Nesse período, adquiri experiência com arquitetura de software, computação em nuvem, pipelines de CI/CD e processos de desenvolvimento que mudaram completamente minha forma de enxergar um projeto de software. Assumir a função de AGES III deixou de ser uma preocupação e passou a ser uma oportunidade de aplicar, na prática, conhecimentos que até então eu utilizava apenas em outros contextos profissionais. Mais do que desenvolver funcionalidades, pude participar ativamente da definição da arquitetura da solução, da construção da infraestrutura, da criação dos processos de qualidade e do acompanhamento técnico do time.\par

% text-p-1313
\indent Outro aspecto muito significativo foi perceber que a atuação de um AGES III vai além da implementação de código. Ao longo do projeto, tive a oportunidade de orientar colegas, revisar entregas, discutir soluções técnicas e contribuir para que o time evoluísse como um todo. Em diversos momentos, aprender e ensinar aconteceram simultaneamente, especialmente nas atividades desenvolvidas em conjunto com colegas que estavam iniciando sua trajetória na graduação. Essa troca de conhecimento foi, sem dúvida, uma das experiências mais gratificantes desta edição.\par

% text-p-1314
\indent Por fim, encerro esta participação levando não apenas os conhecimentos técnicos adquiridos, mas também diversos aprendizados sobre liderança técnica, organização de equipes e condução de projetos de software. As experiências vividas como AGES III me proporcionaram uma visão muito mais ampla sobre o desenvolvimento de software e serviram de preparação para o próximo desafio dentro da disciplina: atuar como AGES IV. Levo comigo exemplos de boas práticas, situações que podem ser aprimoradas e a convicção de que cada edição da AGES contribui para formar profissionais mais preparados para os desafios encontrados no mercado de trabalho.\par

